\begin{abstract}
Layer skipping makes Transformer inference cheaper by deleting entire residual
blocks, but the operator it uses to do so is crude: it assumes the missing
residual update is exactly zero. We ask whether that update is instead
\emph{predictable} from computation the model has already performed, and show
that it largely is. Consecutive residual updates trace a locally smooth
trajectory across depth, so a block's update can be forecast from the updates
that precede it. We introduce \textbf{\mname (Depth-wise AutoRegressive update
prediction)}, which replaces a skipped block with a short autoregressive forecast
of its residual update, using one or two scalars per skipped layer fitted in
closed form on a handful of unlabeled calibration sequences. \mname needs no
gradient training, adds only $O(Td)$ scaled vector additions, and leaves hidden
width and attention structure untouched. Across Qwen2.5 models from 0.5B to
\topscale parameters, at identical skipped-layer sets, \mname recovers \ph{62}\%
of the quality that plain skipping discards: with four skipped blocks it costs
\ph{1.3} points of average downstream accuracy where plain skipping costs
\ph{3.4}, while retaining the same block-compute reduction and matching plain
skipping's prefill latency to within \ph{2}\%. Plain skipping does not merely
remove computation --- it discards computation the model has already forecast.
\end{abstract}
